\chapter{Conclusiones}
\label{ch:conclusiones}

Esta tesis se propuso construir desde cero un motor \textit{ab initio} de estructura atómica con base de B-splines, y emplearlo para seguir la ruptura del acoplamiento de Russell-Saunders a lo largo de la secuencia $np^2$. Los dos objetivos se cumplieron, y el segundo produjo además un conjunto de discrepancias con el experimento que resultaron ser más informativas que el acuerdo mismo. Este capítulo recoge lo que el trabajo estableció, delimita lo que quedó fuera de su alcance y señala las extensiones que cada limitación define de manera natural.

\section{Resultados Principales}

El primer resultado es que la base de B-splines resuelve el límite Hartree-Fock no relativista de toda la secuencia sin cambios de arquitectura. Del carbono al estaño, con la carga nuclear multiplicándose por más de ocho y el número de electrones por más de cuatro, bastó ajustar la densidad de nudos y el parámetro de la malla exponencial para acomodar el rango radial creciente. Las energías totales coinciden con las de Froese Fischer dentro de $3.3\times10^{-5}$ Hartree, las integrales de Slater dentro de $1.8\times10^{-5}$, y el cociente del virial se aparta de $2$ en menos de $1.5\times10^{-5}$ en los cuatro elementos, según el Cuadro \ref{tab:resultados_energia_global}. Esa precisión es la que permite atribuir las discrepancias espectroscópicas posteriores a la física del modelo y no a su implementación numérica.

El segundo resultado es de carácter algorítmico y fue condición necesaria para el primero. El ciclo de punto fijo no converge en el germanio ni en el estaño, y el desplazamiento de nivel que lo estabiliza lo vuelve impracticablemente lento. La extrapolación de Pulay sobre el conmutador proyectado descrita en la Sección \ref{sec:cdiis_metodo} reduce el conteo de iteraciones por factores de $2.8$ y $4.8$ en esos dos elementos, y lo hace sin desplazar el punto fijo: las energías obtenidas con y sin aceleración coinciden por debajo de $10^{-8}$ Hartree. El diagnóstico que hizo falta para llegar ahí tiene interés propio, ya que el vector de error de Pulay en su forma habitual no se anula en la solución de este esquema ROHF, y hubo que proyectarlo fuera del espacio ocupado para recuperar una medida legítima del error.

El tercer resultado es el físico. La secuencia documenta de manera cuantitativa el tránsito entre regímenes de acoplamiento, medido tanto por la razón de intervalos del término fundamental, que cae de $1.98$ a $1.14$ en el cálculo y de $1.64$ a $1.03$ en el experimento, como por el peso del carácter $^1D_2$ en el nivel $^3P_2$, que pasa de ser indetectable en el carbono a casi un diez por ciento en el estaño. Con la polarización del core calibrada, los intervalos del tripleto del germanio y del estaño se reproducen dentro del dos por ciento, y el modelo captura correctamente el ordenamiento de los cinco niveles en los cuatro elementos.

El cuarto resultado es metodológico y quizá el más útil para quien continúe este trabajo. Cada discrepancia relevante con el experimento pudo asociarse a una aproximación concreta del modelo, y en cada caso la asociación predice correctamente el signo del error. Un parámetro de espín-órbita de un cuerpo que sobra en el carbono y falta en los pesados, un core congelado que deja alto al término $^1S$ y una polarizabilidad ajustada que corrige el desdoblamiento magnético a costa de la energía de ionización no son fallas independientes, sino tres manifestaciones de las tres aproximaciones que definen el modelo. Establecer esa correspondencia, y verificarla numéricamente, resultó más instructivo que perseguir un acuerdo global mediante parámetros adicionales.

\section{Alcance del Trabajo}

Conviene delimitar con precisión lo que el modelo no describe, porque marca la frontera entre lo que puede mejorarse refinando el cálculo y lo que exige un formalismo distinto. El Capítulo \ref{ch:discusion} desarrolla cada punto; aquí se enuncian en su forma más compacta.

El Hamiltoniano de estructura fina se reduce a un único operador de un cuerpo parametrizado por $\zeta_{np}$, de modo que los términos de dos cuerpos de Breit-Pauli quedan fuera. Su ausencia es visible en el carbono, donde el modelo cumple la regla de intervalos de Landé y el átomo real no. El espacio de correlación se limita a la pareja de valencia sobre un core congelado, lo que excluye la casi degeneración entre las configuraciones $ns^2np^2$ y $np^4$ y deja alto al término $^1S$. Los orbitales son solución de un problema no relativista, de manera que la contracción de origen relativista, que en los elementos pesados el experimento reclama, solo puede entrar por la vía indirecta del parámetro $\alpha_d$. Y tres decisiones de implementación, la ortogonalización por Gram-Schmidt, la energía extraída de autovalores y el truncamiento en $l_{\max}=3$, afectan la precisión en magnitudes uno o dos órdenes por debajo de los errores físicos anteriores.

\section{Trabajo Futuro}

Las limitaciones anteriores no son igualmente costosas de levantar, y conviene ordenarlas por el esfuerzo que exigen frente a lo que devolverían.

La extensión de menor costo es completar el Hamiltoniano de Breit-Pauli con sus términos de dos cuerpos, en particular el de espín-espín. No requiere tocar el ciclo autoconsistente ni el cálculo de correlación, solo añadir elementos de matriz sobre la misma base de determinantes de valencia y con las mismas integrales radiales. Su efecto esperado es corregir la forma del multiplete en los elementos ligeros, donde la razón de intervalos medida se aparta de la de Landé en un dieciocho por ciento sin que el modelo actual pueda reproducir esa desviación.

De costo intermedio es sustituir la ortogonalización por Gram-Schmidt por multiplicadores de Lagrange dentro del problema variacional. Eliminaría la desviación de entre el $0.2\%$ y el $0.3\%$ en el momento $\langle r^{-3} \rangle$ de los elementos con core $p$, y de paso haría que el vector de error de Pulay se anulara en la solución sin necesidad de proyección. El precio es que obligaría a rehacer todos los cálculos autoconsistentes y, con ellos, los de interacción de configuraciones que se apoyan en sus orbitales.

La extensión de mayor alcance, y la única que resolvería el error de los singletes, es un cálculo multirreferencial de cuatro electrones con las subcapas $ns$ y $np$ simultáneamente activas. La Sección \ref{ch:discusion} muestra que no puede plantearse como una corrección añadida al esquema actual, ya que los dos mecanismos describen en parte la misma física y sumarlos sobrecorrige el término $^1S$ por un factor cercano a cuatro en el carbono. Se trata, por tanto, de un motor distinto y no de una mejora incremental, y esa es la razón por la que quedó fuera del alcance de esta tesis.

En paralelo a lo anterior, dos líneas permitirían separar las contribuciones que hoy absorbe el parámetro $\alpha_d$. La primera es incorporar el término de dos cuerpos del potencial de polarización deducido en la Ecuación \eqref{eq:cpp_dos_cuerpos}, que describe el apantallamiento dieléctrico que el core interpone entre los dos electrones de valencia y que reduciría la repulsión efectiva entre ellos. Es un candidato natural para corregir el término $^1D_2$, que queda alto en los cuatro elementos y que ninguna corrección de la parte radial ha logrado bajar. La segunda es obtener polarizabilidades dipolares de los cores del germanio y del estaño por un cálculo independiente sobre los iones de capa cerrada correspondientes, lo que permitiría contrastar el valor ajustado contra uno físico y cuantificar cuánta de la corrección es polarización genuina y cuánta es relatividad disfrazada.

Finalmente, la ruta que cerraría el problema en su raíz es abandonar el punto de partida no relativista y resolver las ecuaciones de Dirac-Hartree-Fock. El déficit de $\zeta_{np}$ documentado en la Sección \ref{sec:zeta_diagnostico}, que en el germanio y el estaño supera lo que produciría un núcleo desnudo sobre los orbitales calculados, apunta a que la función de onda debería estar más contraída de lo que un tratamiento de Schrödinger permite. Existe además precedente directo de que la base empleada aquí es adecuada para ese fin, ya que los B-splines han demostrado ser estables en la resolución del problema relativista para elementos pesados \cite{Johnson1988}, lo que sugiere que la arquitectura construida en esta tesis admite esa extensión sin rehacerse por completo.

\section{Consideración Final}

El propósito declarado de este trabajo en la Sección \ref{sec:planteamiento} no era aportar un desarrollo teórico original sino demostrar dominio del problema multielectrónico mediante la construcción completa de un motor de cálculo. Vale la pena señalar en qué sentido ese propósito se cumplió más allá de la validación numérica. Un código que reproduce valores tabulados demuestra que las ecuaciones se resolvieron correctamente; lo que demuestra comprensión es poder explicar por qué los valores que no reproduce fallan en la dirección en que fallan. Los diagnósticos del Capítulo \ref{ch:discusion}, que acotan el parámetro de espín-órbita contra la cota de núcleo desnudo, que aíslan la casi degeneración midiendo su tamaño sin implementarla, y que distinguen la desviación de Landé del carbono de la del estaño pese a que ambas apuntan en el mismo sentido, son el producto que este trabajo entrega junto con el código. La secuencia $np^2$ resultó un banco de pruebas afortunado precisamente por eso: cuatro elementos con la misma estructura angular bastan para que cada aproximación del modelo se vuelva visible en un elemento distinto.

%%% Local Variables:
%%% mode: LaTeX
%%% TeX-master: "../main"
%%% End:
